La gestión de los proyectos de extensión en la Universidad del Valle de Guatemala (UVG) representaba un desafío
operativo y técnico significativo. La falta de un sistema centralizado debido a la dependencia histórica de
procesos manuales y canales de comunicación descentralizados, y en ocasiones informales, resultaba en una gestión
ineficiente y propensa a errores humanos.
Por lo que surgió la necesidad de modernizar este proceso mediante la implementación de un sistema de gestión de
proyectos de extensión por medio de una plataforma web. Al evaluar los resultados obtenidos tras la implementación
de la plataforma, se evidencia que la adopción de un nuevo modelo digital unificado no solo es técnicamente viable,
sino que también es esencial para mejorar la eficiencia operativa y la calidad de los servicios ofrecidos por la
UVG. Esta plataforma permite responder de manera directa a las necesidads de gobernanza y trazabilidad de los
proyectos de extensión.

La API, producto del desarrollo del \textit{Backend}, se diseñó con una arquitectura modular y escalable. Las pruebas
de estrés realizadas demostraron que la API puede manejar un alto volumen de peticiones concurrentes sin comprometer
el rendimiento, lo que es crucial para garantizar la disponibilidad y capacidad de respuesta de la plataforma. Las
pruebas de estrés incluyeron tres escenarios de carga: baja, media y alta, con 50, 100 y 250 usuarios simultáneos
respectivamente.
En promedio, la API mantuvo tiempos de respuesta inferiores a 200ms. Con un tiempo de respuesta promedio de 16 a 21
milisegundos para consultas de datos como se puede observar en la Figura \ref{fig:stressTestProjects}, y 116 a 283
milisegundos para operaciones que involucran escritura de datos o procesos más complejos como procesamiento criptográfico,
como el inicio de sesión en la Figura \ref{fig:stressTestAuthLogin}.
Este comportamiento demuestra que la API fue capaz de manejar alta demanda concurrente sin degradar significativamente
el rendimiento, lo que es un indicador positivo de su escalabilidad y robustez. Esto gracias a la elección de tecnologías
modernas y eficientes, como Node.js y Express, que permiten manejar múltiples conexiones simultáneas de manera eficiente
por su modelo de ejecución basado en eventos de entrada/salida no bloqueante, soportando picos de tráfico sin comprometer
la capacidad de respuesta.
Además, el diseño modular, integración de pruebas unitarias y pruebas de funcionalidad facilitan la incorporación
de nuevas funcionalidades o bien, la actualización de las existentes. Siendo capaz de adaptarse a las necesidades
cambiantes de la UVG, lo que es fundamental para la mantenibilidad y sostenibilidad a largo plazo del \textit{Backend}.

La API no fue diseñada únicamente para ser modular y escalable, sino que también segura. Los \textit{endpoints}
fueron diseñados para ser accesibles solo a usuarios autenticados y autorizados, utilizando JWT para la autenticación
y control de acceso basado en roles. Las pruebas de funcionalidad confirmaron que el diseño de los \textit{endpoints} 
protegidos y no protegidos por autenticación y autorización se implementó correctamente, como se puede observar en
la Figura \ref{fig:updateProject}, donde únicamente el dueño del proyecto puede actualizar la información del proyecto,
mientras que otros usuarios no autorizados reciben un error de acceso denegado. Al igual que en la Figura
\ref{fig:getProjectsUnauthorized}, donde un usuario no autenticado recibe un error de acceso denegado al intentar acceder 
lista de proyectos de extensión. El control de acceso basado en roles impone reglas de negocio estrictas necesarias
para salvaguardar la integridad de los datos y garantizar que únicamente los usuarios autorizados puedan realizar
acciones específicas, lo que es fundamental para proteger la información sensible y mantener la confianza de los
usuarios en la plataforma.
De igual manera, la solidez de la seguridad de la API se ve reforzada por el análisis de vulnerabilidades automático
con OWASP ZAP, que logró identificar alertas de seguridad medias y bajas, pero no encontró vulnerabilidades críticas.
Esto garantiza que la API proporciona un entorno seguro para los usuarios de la plataforma, al estar protegida contra
amenazas comunes. Por supuesto, las alertas de seguridad medias y bajas identificadas fueron abordadas y mitigadas en
la medida de lo posible para mejorar aún más la seguridad general de la API. Como la sanitización de entradas para
prevenir ataques de inyección de código malicioso, entre otros.

El \textit{Frontend} se implementó utilizando React, lo que permitió crear una interfaz de usuario moderna, interactiva
y \textit{responsive}.
La selección de una arquitectura web \textit{Single Page Application} (SPA, por sus siglas en inglés) utilizando la estrategia
de renderizado web \textit{Client-Side Rendering} (CSR, por sus siglas en inglés) fue lo que permitió al usuario interactuar
con la plataforma de manera fluida, sin interrupciones causadas por recargas de página, lo que mejora significativamente
la experiencia.
Las tecnologías seleccionadas para el desarrollo del \textit{Frontend} permitieron crear un diseño fiel a los prototipos
aprobados y una experiencia de usuario fluida. TypeScript, EsLint, Prettier y Storybook contribuyeron significativamente
a la calidad, consistencia y mantenibilidad del código, lo que se traduce en una interfaz de usuario más robusta y fácil
de mantener a largo plazo. Material UI se utilizó para garantizar la fidelidad visual con los prototipos aprobados, lo que
a su vez contribuyó a una experiencia de usuario coherente y atractiva. TanStack Query fue seleccionado para la gestión de
datos provenientes de la API, permitió una gestión eficiente del estado de los datos, mejorando el rendimiento al tener
control sobre la caché y las actualizaciones de datos. TanStack Forms fue utilizado para la gestión de formularios, lo que
permitió validaciones robustas en tiempo real, mejorando la experiencia del usuario al proporcionar retroalimentación
inmediata sobre la validez de los datos ingresados y deshabilitando el envío de formularios con datos inválidos,
lo que contribuyó a la integridad de los datos y a una experiencia de usuario más amigable.

Por lo que, todo lo anterior contribuyó a los resultados obtenidos en las pruebas de usabilidad por parte de usuarios
finales, quienes en general reportaron una experiencia muy positiva con la plataforma.
Como se observó en la Figura \ref{fig:completedTasksGraph}, todos los usuarios fueron capaces de completar las tareas
asignadas. Aunque es importante destacar que algunos usuarios interpretaron mal la escala de dificultad, por lo que
colocaron tareas que consideraron fáciles como difíciles; sin embargo, la mayoría de los participantes reportaron que
las tarea fueron muy fáciles o fáciles de completar, lo que indica que la plataforma fue intuitiva. Además, en la Figura
\ref{fig:difficultyAndTimeGraph} se observa que, en promedio, la mayoría de usuarios completaron las tareas en menos de
3 minutos, lo que refuerza la percepción de que la plataforma requiere una curva de aprendizaje baja.
El \textit{Frontend} no solo ofreció una experiencia agradable para los usuarios, sino que también se integró de manera
efectiva con el \textit{Backend}, lo que permitió realizar las operaciones requeridas que comprenden la gestión de
proyectos de extensión, incluyendo otras operaciones como la gestión de usuarios, autenticación, entre otros. Esto se
puede confirmar con las pruebas de funcionalidad realizadas, donde luego de completar las tareas asignadas, se verificó
que el estado de la plataforma reflejaba correctamente las acciones realizadas por los usuarios, tal como se puede observar
en las Figuras \ref{fig:usabilityTestDB} y \ref{fig:usabilityTestBucketS3}, donde la Base de Datos y el Bucket S3 respectivamente,
reflejan correctamente los cambios realizados por un participante durante las pruebas de usabilidad.
Además, las pruebas de usabilidad permitieron obtener el puntaje promedio de la Escala de Usabilidad del Sistema (SUS, por sus siglas
en inglés), que fue de 80. Este puntaje, dentro de la escala de SUS, se considera un puntaje excelente. Esto indica que los
usuarios finales percibieron la plataforma como altamente usable, intuitiva y genera un alto nivel de satisfacción, lo que a su
vez refuerza la aceptación y adopción de la plataforma por parte de los usuarios finales como el nuevo sistema de gestión de
proyectos de extensión en la UVG.
Juntamente con el \textit{Backend}, el \textit{Frontend} fue sometido al análisis de vulnerabilidades automático con OWASP ZAP,
el cuál identificó alertas de seguridad medias y bajas, las cuales fueron abordadas y mitigadas en la medida de lo posible como
la definición de políticas de seguridad de contenido (CSP, por sus siglas en inglés) para proteger contra ataques de inyección
de código malicioso, como ataques de Cross-Site Scripting (XSS, por sus siglas en inglés) comunes.

Es importante destacar que las pruebas de usabilidad también revelaron áreas de mejora, como la necesidad de mover ciertos elementos
visuales como el "Registro de horas" por parte de estudiantes, a una ubicación más visible dentro de la plataforma, no dentro del
historial de proyectos finalizados, ya que no para todos los usuarios resultó intuitivo. O bien la necesidad de agregar un tutorial
inicial para guiar a los usuarios en el uso de la plataforma. Y otros comentarios en relación al diseño visual. Los cuales deberán
de ser considerados para futuras iteraciones de la plataforma.

Resulta permitente afirmar, entonces, que la implementación de la plataforma web para la gestión de proyectos de extensión en la
UVG logró cumplir con los objetivos planteados, al desarrollar un sistema que centraliza la administración, publicación, participación
y seguimiento de los proyectos de extensión, permitiendo a estudiantes, directores y organizaciones externas participar oportunamente
en el ciclo de vida de los proyectos de extensión. 